\documentclass[a4paper,11pt]{article}

% page geometry
\usepackage[colorlinks,linkcolor=blue,citecolor=blue,urlcolor=blue]{hyperref}
\usepackage[lmargin=2cm,rmargin=2cm,tmargin=1.25cm,bmargin=1.25cm,includefoot,includehead]{geometry}

% typography
\usepackage[T1]{fontenc}
\usepackage[utf8]{inputenc}
\renewcommand{\rmdefault}{ptm}   % Times-like serif
\renewcommand{\sfdefault}{phv}   % Helvetica-like sans serif

% text structures
\usepackage{enumitem}
\usepackage{tabu}
\usepackage[dvipsnames]{xcolor}
\usepackage{tcolorbox}
\usepackage{colortbl}
\usepackage{graphicx}
\usepackage{xspace}
\usepackage{amsmath}
\usepackage{booktabs}
\usepackage{array}
\usepackage{fancyvrb}

% colors
\definecolor{ANRblue}{HTML}{004d95}
\definecolor{ANRred}{HTML}{FF0000}
\definecolor{ANRpurple}{HTML}{800080}

% useful packages
\usepackage{siunitx}
\usepackage[version=3]{mhchem}
\usepackage{csquotes}

% comments / draft commands
\newcommand{\Mathieu}[1]{\textcolor{orange}{#1}}
\newcommand{\Helene}[1]{\textcolor{purple}{#1}}

\newcolumntype{P}[1]{>{\raggedright\arraybackslash}p{#1}}

% section title format
\usepackage{tocloft}
\renewcommand{\cfttoctitlefont}{\color{ANRblue}\normalfont\scshape\bfseries\sffamily\Large}

\usepackage{titlesec}
\titleformat{\section}
{\color{ANRpurple}\normalfont\scshape\bfseries\sffamily\Large}
{\color{ANRpurple}\thesection}{1em}{}
\renewcommand{\cftsecfont}{\color{ANRpurple}\normalfont\scshape\bfseries\sffamily\large}

\titleformat{\subsection}
{\color{ANRblue}\normalfont\scshape\bfseries\sffamily\large}
{\color{ANRblue}\thesubsection}{1em}{}
\renewcommand{\cftsubsecfont}
{\color{ANRblue}\normalfont\scshape\bfseries\sffamily}

\titleformat{\subsubsection}
{\color{ANRpurple}\normalfont\bfseries\sffamily}
{\color{ANRpurple}\thesubsubsection}{1em}{}
\renewcommand{\cftsubsubsecfont}
{\color{ANRpurple}\normalfont\bfseries\sffamily\small}

\titleformat{\paragraph}[runin]
{\color{ANRblue}\normalfont\sffamily}
{\color{ANRblue}\theparagraph}{1em}{}

% project info
\newcommand{\mytitle}[0]{Bayesian Modeling of Emotional Dynamics Using Hidden Markov Models}
\newcommand{\mycourse}[0]{STAT405 Bayesian Statistics Course Project}
\newcommand{\myterm}[0]{Project Proposal}

\author{
Marc Shi (56215593) \\
Yu Chang (47945050) \\
\small \texttt{lianjies@student.ubc.ca} \\
\small \texttt{changy72@student.ubc.ca}\\
\small Department of Statistics \\
The University of British Columbia \\
}

% headers
\usepackage{fancyhdr}
\pagestyle{fancy}
\lhead{\mycourse}
\rhead{\Large\textcolor{ANRblue}{\myterm}}
\renewcommand{\headrulewidth}{0pt}

\lfoot{}
\cfoot{}
\rfoot{\thepage}

\title{\vspace{-\baselineskip}\mytitle}
\date{\today}

\newcommand{\content}[1]{\emph{#1}\\}
\begin{document}
\maketitle
\thispagestyle{fancy}
\section{Project repository}
The public repository for this project, containing the dataset, code, and model implementations, is available at:
\url{https://github.com/Marclianjie/STAT405-Project-/tree/main}.
The repository will be regularly updated throughout the ongoing of the project and will contain all scripts used for data processing, model implementation, and posterior inference.

\section{Data}Our interest in this project is to model emotional states using repeated measurements and to study how these states evolve over time. Based on this research theme, we will consider one primary dataset and one backup dataset for the analysis:
\begin{itemize}
   \item Our \textbf{main} dataset consists of Positive and Negative Affect Schedule (PANAS)-style affect items adapted for the Experience Sampling Method (ESM), associated with the study by Rowland and Wenzel \cite{Rowland2020}. This scale contains two main constructs: positive affect (happy, enthusiastic, active, alert) and negative affect (sad, anxious, irritated, tired), with each emotional state measured by a single item.

The dataset contains 95 participants with repeated emotional measurements collected approximately 10 times per day over 6 consecutive days. The dataset is publicly available \href{https://osf.io/jmz2n/files/m5fcy}{here} and the original studies are under the \href{https://github.com/jmbh/EmotionTimeSeries}{EmotionTimeSeries Repository}. Below we show an example of the head of the dataset.

\begin{verbatim}
subjno dayno beep group emo1_m emo2_m emo3_m emo4_m emo5_m emo6_m emo7_m emo8_m
1      1     1    2     79     64     30     74     19     7      8      58
1      1     3    2     20     24     37     36     62     41     73     67
1      1     4    2     25     24     33     29     42     68     77     81
1      1     5    2     37     24     40     17     35     36     12     76
1      2     2    2     64     39     36     41     36     28     35     59
\end{verbatim}
with \texttt{emo1\_m} to \texttt{emo8\_m} representing the item scores (ranging from 0 to 100) for the eight emotion states corresponding to positive and negative affect.

\textbf{Difference from prior work}: Rowland and Wenzel \cite{Rowland2020} used the longitudinal data to estimate affect networks based on partial correlations between emotions, without explicitly modeling the temporal evolution of emotional states. In contrast, our approach treats the data as a time series and uses a Bayesian Hidden Markov Model to infer latent emotional states and their transition dynamics.

\item Our \textbf{backup} dataset is drawn from the \href{https://github.com/mkullar/DataDrivenEmotionDynamics/tree/main}{DataDrivenEmotionDynamics} repository associated with the study by \href{https://psycnet.apa.org/fulltext/2023-75204-001.html}{Kullar et al} \cite{Kullar2024}. It is an experience sampling dataset with 105 participants and repeated within-person emotional measurements over time, making it suitable for the same modeling framework as our main dataset. The dataset records repeated assessments of multiple affective states on a 1 - 7 scale, along with emotion duration and mind-wandering information. 
 
Because each observation corresponds to a time point within an individual, the data can be treated as a multivariate longitudinal time series and used to fit a Bayesian Hidden Markov Model for latent emotional states and their temporal transitions. Below we show an example of the head of the backup dataset.

\\

{\scriptsize
\begin{verbatim}
id day beep angry emotion_chronometry enthusiastic happy irritated mind_wandering nervous pleased relaxed sad stressed
1  1   2    1     1                    5            6     6         0              1       5       7       1   1
1  1   3    1     1                    5            5     3         0              1       6       7       1   1
1  1   4    1     60                   4            5     6         1              1       4       4       1   2
1  1   5    1     1                    4            6     1         0              1       4       7       1   1
1  2   2    1     30                   6            6     1         1              1       7       6       1   1
1  2   3    1     1                    6            6     1         1              1       6       7       1   1
\end{verbatim}
}


The variables shown in the dataset head include \texttt{id}, \texttt{day}, and \texttt{beep}, which identify the participant, study day, and within-day measurement occasion, respectively. The variables \texttt{angry}, \texttt{enthusiastic}, \texttt{happy}, \texttt{irritated}, \texttt{nervous}, \texttt{pleased}, \texttt{relaxed}, \texttt{sad}, and \texttt{stressed} represent self-reported emotion intensities at each measurement occasion. The variable \texttt{emotion\_chronometry} records the reported duration of the current emotional experience, and \texttt{mind\_wandering} is an indicator of whether mind wandering occurred at that time point. Together, these variables provide repeated multivariate emotion profiles that are well suited for modeling latent affective states.

\textbf{Why is this treated as Backup}: Although it has a similar repeated-measures structure, it is somewhat less directly aligned with our primary modelling goal. Compared with the main dataset, its emotion variables are measured on a coarser 1-7 scale and appear to have less clean observation structure.



\end{itemize}
% Main dataset, backup dataset, structure of dataframe, research question.
\section{Research question}

The main goal of this project is to investigate the temporal dynamics of emotional states using Bayesian modeling. In particular, we aim to address the following questions:

\begin{enumerate}
    \item To what extent does a Bayesian Hidden Markov Model provide a better fit to the observed emotional data compared to a naive model that assumes independence across time? We will assess this by comparing predictive performance and posterior predictive errors.

    \item Can the Bayesian Hidden Markov Model be used to infer and predict the latent emotional state of a participant at the next measurement occasion?

    \item How do different Bayesian posterior inference methods compare in terms of computational efficiency and convergence behavior when applied to the Hidden Markov Model?
\end{enumerate}
\section{Proposed methodology \& comparison plan}

In this project, we will analyze the EMA emotion dataset using Bayesian methods. 
Our analysis will involve two Bayesian models and two posterior computation methods, 
allowing us to compare both modeling approaches and inference algorithms.

\subsection{Bayesian models}

We will implement and compare the following two models:

\begin{enumerate}

\item \textbf{Naive model.}  
The naive model assumes that emotional states are independent across time. 
For each measurement occasion, the latent true emotional state is assumed to be an i.i.d.\ draw from a prior distribution. 
The observed emotion score is then generated from a likelihood distribution conditional on this latent state. 
Under this model, both the latent emotional states and the observed measurements are assumed to be independent across time, 
therefore ignoring any temporal dependence in the longitudinal data. 
This model serves as a baseline that does not incorporate the time-series structure of the dataset.

\item \textbf{Hidden Markov Model (HMM).}  
The Hidden Markov Model builds upon the naive model by assuming the latent emotional states evolve over time 
according to a Markov transition process, in which one emotional state of time t is based on the previous emotional state of time t-1. And the observed emotion scores are assumed to be 
generated from distributions that depend on the underlying latent state. This model 
allows us to capture temporal dependence and infer the sequence of hidden emotional 
states that give rise to the observed affect measurements.

\end{enumerate}

\subsection{Posterior computation methods}

For posterior inference, we will apply two Bayesian computation techniques introduced 
in this course:

\begin{enumerate}

\item \textbf{Self-Importance Sampling (SNIS).}  
This method approximates the posterior distribution by drawing samples from a proposal 
distribution and reweighting them according to their likelihood.

\item \textbf{Markov Chain Monte Carlo (MCMC) using the Metropolis--Hastings algorithm.}  
This method constructs a Markov chain whose stationary distribution is the posterior 
distribution of interest, allowing us to generate samples that approximate the posterior.

\end{enumerate}

\subsection{Model and inference comparison}

We will compare the naive model and the Hidden Markov Model to evaluate the extent to 
which modeling latent emotional states improves the fit to the observed data. Model 
performance will be assessed using predictive accuracy and posterior predictive errors.

In addition, we will compare the two posterior computation methods in terms of 
computational efficiency and convergence behavior when applied to the Hidden Markov 
Model. This comparison will help illustrate the practical differences between sampling 
approaches and importance sampling methods for Bayesian inference.


\section{Team Contribution Plan}

The responsibilities for this project will be divided between the two team members to ensure balanced contributions.

Yu Chang will be primarily responsible for the literature review, as well as drafting the introduction and background sections of the report. Marc Shi will be responsible for the methodology implementation and the results analysis sections.

For the coding component, both team members will contribute independently by implementing and running one of the two models considered in this project. In particular, each member will implement and run one model, allowing us to independently verify the results and compare the models more effectively. Marc Shi will compile and present the final results in the report.

Throughout the project, both team members will communicate regularly to coordinate progress, discuss modeling decisions, and ensure the consistency and quality of the final report.

\bibliographystyle{plain}
\clearpage
\bibliography{references}
\end{document}